\section{Robots Manipuladores en Entornos Académicos}

El uso de robots manipuladores en entornos académicos se ha incrementado significativamente en los últimos años gracias a la disponibilidad de plataformas abiertas y de bajo costo. Entre estas plataformas, el PhantomX Pincher destaca como una opción ampliamente utilizada en laboratorios universitarios debido a su arquitectura modular, compatibilidad con los servomotores Dynamixel y facilidad de integración con sistemas de control como ROS (Robot Operating System).

La comunidad de desarrollo ha contribuido significativamente a la expansión de las capacidades de esta plataforma, con repositorios como el de \textit{snt-spacer} \cite{sntspacer}, que proporciona paquetes de software tanto para ROS 1 como para ROS 2, facilitando la integración del PhantomX Pincher en proyectos educativos y de investigación modernos.

\section{Kits Educativos Comerciales}

A nivel internacional, empresas como Elephant Robotics han desarrollado kits educativos como el myCobot AI Kit \cite{ElephantRobotics2021}, que integran manipulación robótica, visión por computador y módulos de inteligencia artificial en entornos compactos y accesibles. Estos kits han sido diseñados para ser utilizados en cursos de introducción a la robótica, visión artificial y programación, y se han convertido en referentes comerciales para la creación de entornos didácticos orientados a la enseñanza práctica de tecnologías emergentes.

La evolución de estos kits ha sido notable, con actualizaciones continuas que incorporan las últimas tendencias en robótica educativa e investigación \cite{ElephantRobotics2023}, demostrando la importancia de mantener plataformas de enseñanza actualizadas con las tecnologías emergentes.

\section{ROS2 y Sistemas Embebidos}

En el ámbito académico, múltiples universidades y centros de investigación han explorado la combinación de ROS con sistemas de visión por computador como OpenCV, junto con plataformas embebidas como Raspberry Pi, para desarrollar proyectos de clasificación, seguimiento de objetos o manipulación autónoma.

Erős et al. \cite{Eros2019} han demostrado la viabilidad de arquitecturas de comunicación basadas en ROS2 para sistemas de automatización colaborativa e inteligente, destacando las ventajas de la nueva generación del sistema operativo robótico en términos de tiempo real, seguridad y escalabilidad.

Por su parte, Meijer \cite{Meijer2021} ha investigado específicamente el uso de Raspberry Pi con ROS2, implementando frameworks de software en tiempo real utilizando Xenomai. Este trabajo demuestra que es posible obtener rendimiento determinístico en plataformas embebidas de bajo costo, lo cual es fundamental para aplicaciones de control robótico que requieren respuestas predecibles y oportunas.

Adicionalmente, proyectos de código abierto como el desarrollado por \textit{noshluk2} \cite{noshluk2} han demostrado la integración exitosa de ROS2, Raspberry Pi y sistemas de visión inteligente para robótica, proporcionando ejemplos prácticos de implementación que pueden servir como referencia para desarrollos educativos.

\section{Visión por Computador en Robótica}

Estas experiencias han demostrado la viabilidad de construir sistemas funcionales de bajo costo con capacidades avanzadas de percepción y control, utilizando herramientas de software libre y hardware de fácil acceso.

Ruiz-del-Solar et al. \cite{RuizDelSolar2018} han realizado una revisión comprehensiva sobre métodos de aprendizaje profundo para visión robótica, destacando cómo estas técnicas han revolucionado la percepción artificial en robótica, permitiendo a los sistemas reconocer y clasificar objetos con una precisión sin precedentes.

En el contexto específico de la manipulación robótica, Du et al. \cite{Du2019} han documentado los avances en sistemas de agarre basados en visión, cubriendo desde la localización de objetos hasta la estimación de poses y el cálculo de puntos de agarre para pinzas paralelas. Esta revisión proporciona una perspectiva integral de los desafíos y soluciones en el campo de la manipulación guiada por visión.

\section{Brecha Identificada}

Sin embargo, en el contexto específico del robot PhantomX Pincher, no se han documentado ampliamente kits educativos que integren de forma sistemática elementos como visión artificial, ROS2 y procesamiento embebido con Raspberry Pi. Esta ausencia representa una oportunidad para proponer un diseño conceptual que adapte estas tecnologías a dicha plataforma, tomando como inspiración kits comerciales existentes pero considerando las particularidades técnicas y pedagógicas del entorno universitario local.
